\documentclass[10pt,a4paper]{amsart}
\usepackage[margin=19mm]{geometry}
\usepackage{amsmath,amssymb,mathtools}
\usepackage[scr=rsfs,cal=euler]{mathalfa}
\usepackage{xfrac}
\usepackage[unicode,hidelinks,bookmarks=false]{hyperref}
\newcommand{\ie}{i.e.\ }
\newcommand{\cf}{cf.\ }
\newcommand{\resp}{respectively\ }
\newcommand{\Subsection}{\subsection}
\providecommand{\nobreakdash}{\nobreak\hbox{-}}
\setlength{\emergencystretch}{2em}
\multlinegap0pt
\usepackage{amssymb}
\usepackage{mathtools}
\newtheorem{lemma}{Lemma}
\title{A Complete Answer to Erdős Problem 690\\ {\large Discovered by the Multiscalar Fields System}}
\author{Shouqiao Wang, Davide Crapis}
\date{Source: arXiv:2605.08542v1 (CC BY 4.0)}
\begin{document}
\maketitle

\noindent\emph{Source and licence.} A Complete Answer to Erdős Problem 690\\ {\large Discovered by the Multiscalar Fields System}. Version arXiv:2605.08542v1 (CC BY 4.0), distributed by arXiv under the Creative Commons Attribution 4.0 International licence, http://creativecommons.org/licenses/by/4.0/, as recorded in the arXiv licence metadata for this exact version and shown on the abstract page https://arxiv.org/abs/2605.08542v1. Full licence notice: \texttt{licenses/arxiv-2605.08542-CC-BY-4.0.txt}.\par\smallskip

\noindent\textit{Editorial scope.} Counted unit: the article's lemma lem:Rbounds (source0.tex lines 206--216). The article's complete original proof of it is delivered with it (source0.tex lines 218--288, one \texttt{proof} environment of 71 lines). No mathematics is written, repaired or re-derived: the statement and the proof are the article's own lines, in the article's own notation, and no retained line is substituted or re-flowed.  No further source-local context is needed: every cross-reference of every retained line resolves inside the two delivered spans.  Nothing was omitted from any retained span, so the package carries no omission record.  No retained line contains a citation, so the wrapper carries no bibliography and no bibliography entry.  No retained line contains a figure environment, an image-inclusion command or drawing code, so no image is delivered and the package declares no asset dependency.  Editorial formatting only, in addition: an AMS \texttt{amsart} preamble that restates the article's own declarations and macros used by the retained lines, namely the environments \texttt{lemma} on separate counters, together with the standard packages those lines need.  The retained environments print in this document as Lemma 1 in order of appearance, whereas the article prints the same items with the numbering of its own full document.  The complete native source of the article as distributed in the arXiv e-print for this exact version is bundled unchanged as \texttt{sources/200158/source0.tex}.
\begin{lemma}[Symmetric-polynomial bounds]\label{lem:Rbounds}
Let \(r\ge 1\), and suppose that \(i\ge r\).  Then \(\delta_r(i)>0\), and
\[
        \frac r{A(p_i)}
        \le
        R_r(i)
        \le
        \frac r{A(p_i)-W_{r-1}} .
\]
The denominator in the upper bound is positive because \(i\ge r\).
\end{lemma}

\begin{proof}
For a fixed subset \(S\subseteq\{1,\ldots,i\}\), the contribution of the event that exactly the primes \(p_j\) with \(j\in S\) divide \(n\) is
\[
        \prod_{j\in S}\frac1{p_j}
        \prod_{j\notin S}\left(1-\frac1{p_j}\right)
        =
        \left(\prod_{j=1}^i\frac{p_j-1}{p_j}\right)
        \prod_{j\in S}\frac1{p_j-1}.
\]
Therefore
\[
        \delta_m(i)
        =
        \left(\prod_{j=1}^i\frac{p_j-1}{p_j}\right)E_m(i),
\]
and hence
\[
        R_r(i)=\frac{E_{r-1}(i)}{E_r(i)}.
\]

Let
\[
        \alpha_i:=A(p_i)=\sum_{j=1}^i w_j .
\]
Expanding \(\alpha_iE_{r-1}(i)\), we obtain
\[
        \alpha_iE_{r-1}(i)=rE_r(i)+\Sigma_i,
\]
where
\[
        \Sigma_i
        :=
        \sum_{\substack{S\subseteq\{1,\ldots,i\}\\ |S|=r-1}}
        \left(\sum_{s\in S}w_s\right)\prod_{s\in S}w_s
        \ge 0.
\]
Since \(w_1\ge w_2\ge w_3\ge\cdots\), every set \(S\) of size \(r-1\) satisfies
\[
        \sum_{s\in S}w_s\le w_1+\cdots+w_{r-1}=W_{r-1},
\]
with the right-hand side interpreted as the empty sum when \(r=1\).
Thus
\[
        0\le \Sigma_i\le W_{r-1}E_{r-1}(i).
\]
Dividing the identity
\[
        \alpha_iE_{r-1}(i)=rE_r(i)+\Sigma_i
\]
by \(E_{r-1}(i)>0\), and writing
\[
        s_i:=\frac{\Sigma_i}{E_{r-1}(i)},
\]
we get
\[
        0\le s_i\le W_{r-1},
        \qquad
        R_r(i)=\frac r{A(p_i)-s_i}.
\]
The lower bound follows from \(s_i\ge 0\).  Since \(i\ge r\),
\[
        A(p_i)-W_{r-1}
        =\sum_{j=r}^{i}w_j
        \ge w_r>0.
\]
Together with \(s_i\le W_{r-1}\), this gives
\[
        A(p_i)-s_i\ge A(p_i)-W_{r-1}>0,
\]
and hence the stated upper bound.
\end{proof}

\end{document}
